\chapter*{General Conclusion and Perspectives}
\phantomsection
\addcontentsline{toc}{chapter}{General Conclusion and Perspectives} %adds to the table of contents 
\label{chap:general-conclusion}

This project addressed the need for comparative evidence that separates measurable
software-stack behaviour from assumptions about vendor ecosystems. We defined a
scenario-based method in which ST serves as the reference, equivalent observable behaviour
is established before comparison, build and measurement conditions are recorded, and
footprint or runtime results are interpreted at the relevant architectural layer. MCU
Workflow Studio and the MCU Benchmark Copilot Agent support this method by structuring
linker-map analysis, project preparation, provenance tracking, fairness review, and
semantic-equivalence checks. Their contribution lies in formalising and assisting the
workflow; the validity of each benchmark still depends on its source artifacts, build
records, and hardware evidence.

The completed ST--GD32 comparison provided footprint and reported runtime observations for
five USB and FreeRTOS scenarios. It revealed a recurring trade-off rather than a uniform
advantage for either ecosystem. STM32 used more flash, with much of the classified
difference associated with HAL components such as RCC, but required substantially less RAM
in the USB scenarios. It also showed a lower reported USB initialisation time, while
steady-state USB and FreeRTOS observations remained close. These results support targeted
investigation of RCC footprint and finer USB feature selection, but they do not establish
that the same pattern applies to other devices, compiler settings, or software stacks.

The ST--TI work provided a different type of evidence. Static analysis characterised the
two low-level driver sets, while twelve matched bare-metal projects produced whole-image
flash and RAM measurements. The aggregate result slightly favoured MSPM33 in flash and
STM32C5 in RAM, with scenario-level outcomes varying according to startup costs, retained
handlers, buffers, and application-owned configuration. The comparison therefore cannot be
reduced to a single ecosystem ranking. The FreeRTOS extension produced the opposite flash
pattern: STM32C5 used 27.9--29.6\,\% less flash in all three scenarios, while MSPM33 retained
a 56--76-byte RAM advantage. Runtime observations were collected, but the 144\,MHz versus
32\,MHz clock difference explains nearly all of the basic-processing gap. Different kernel
versions, configuration options, and preemptive primitives limit the cooperative and
preemptive interpretation, so the current evidence does not establish a scheduler-efficiency
ranking.

Several limitations define the scope of these findings. The evaluated boards, SDK versions,
scenarios, IAR configuration, and linker-map accounting rules form part of each result and
limit its generalisation. Static-analysis findings depend on the selected source sets and
tool configurations, while heuristic classifications or quality indicators are not
substitutes for expert review. Aggregate footprint totals across independent projects are
portfolio summaries, not the memory requirement of one deployable application. Finally, the
supporting tools were designed to improve consistency and traceability, but this report does
not contain a controlled productivity study or a complete cross-environment reproducibility
assessment.

Future work should first repeat the ST--TI runtime phase under the controls identified in
Chapter~\ref{chap:vs-ti}. Each scenario should use equal documented clocks, the same FreeRTOS
version and configuration, equivalent synchronisation primitives, a common measurement
interval, equal sample counts, a warm-up policy, and reported dispersion. Energy measurements
may then be added using the same functional boundaries. The benchmark set can subsequently be
extended to additional middleware and vendor platforms, provided that each new comparison
preserves the same evidence and fairness requirements.

The supporting tools also require further validation. MCU Workflow Studio should be tested
against a larger labelled set of linker maps, with mapping accuracy and repeatability
reported separately, and its parser may be extended beyond IAR EWARM. The MCU Benchmark
Copilot Agent should be evaluated through recorded campaigns that distinguish successful
automation, manual intervention, unresolved unknowns, and hardware-dependent steps. These
extensions would broaden the evidence base without weakening the central rule established
by this work: no comparative conclusion should exceed the measurements and controls that
support it.